\documentclass[11pt]{article}
\usepackage[margin=1in]{geometry}
\usepackage{amsmath,amssymb,amsthm}

\newtheorem{fact}{Fact}
\newtheorem{counterexample}{Counterexample}

\title{Disproof of TLMC Conjecture 00000000492\\
\large A ``random hook length formula'' expectation conjecture}
\author{TLMC submission}
\date{October 2, 2026}

\begin{document}
\maketitle

\begin{abstract}
Conjecture 00000000492 asserts that the expectation of the multiset of hook
lengths of a uniformly random standard tableau, normalized as $(\sum
\mathrm{hooks})/n^2$, equals the closed form $(n^2-1)/3$. We show this is
impossible: the left-hand side is capped at $1$ for every tableau with $n$ cells,
while the conjectured value exceeds $1$ for every $n \ge 3$. A concrete
counterexample ($n=3$, row shape) gives $2/3 \ne 8/3$. Alternative readings
(square shapes, division by $n$) fail as well. Verdict: \textbf{FALSE}.
\end{abstract}

\section{The conjecture}

\begin{quote}
A random version of the hook length formula. \emph{Conjecture:} the expectation
of the multiset of hook lengths of a uniformly random standard tableau equals
$(\sum \mathrm{hooks})/n^2$, explicitly $(n^2-1)/3$.
\end{quote}

We recall that for a Young diagram (partition) $\lambda$, the \emph{hook length}
of a cell $c$ is the number of cells strictly to the right of $c$ in its row,
plus the number of cells strictly below $c$ in its column, plus one. A basic but
decisive observation is that the multiset of hook lengths depends only on the
shape $\lambda$, not on the filling; hence the ``expectation over uniformly
random standard tableaux of the shape'' is a constant attached to $\lambda$
itself, and any universal bound on hook sums of a shape applies to the
expectation verbatim.

\section{The attack: a universal cap that the conjectured value exceeds}

\begin{fact}[Cap]\label{fact:cap}
Let $T$ be any standard tableau with $n$ cells. Then every hook length of $T$ is
at most $n$, because the hook of a cell consists of the cell itself together
with cells strictly right/below, and there are only $n$ cells in total.
Therefore
\[
\sum_{c \in T} h(c) \;\le\; n \cdot n \;=\; n^2,
\qquad\text{so}\qquad
\frac{\sum h(c)}{n^2} \;\le\; 1 .
\]
\end{fact}

\begin{fact}[The conjectured value overshoots the cap]\label{fact:overshoot}
For every $n \ge 3$,
\[
\frac{n^2-1}{3} \;>\; 1,
\]
since $n^2 - 1 > 3 \iff n^2 > 4$.
\end{fact}

Combining Facts~\ref{fact:cap} and~\ref{fact:overshoot}: for every tableau with
$n \ge 3$ cells, of \emph{any} shape,
\[
\frac{\sum h(c)}{n^2} \;\le\; 1 \;<\; \frac{n^2-1}{3},
\]
so the conjecture is false for all $n \ge 3$, regardless of how the
``random tableau'' is distributed. The randomization is irrelevant: the
conjectured value lies strictly above the largest value the left-hand side could
possibly take.

\section{Concrete counterexample}

\begin{counterexample}[$n = 3$, row shape $(3)$]\label{ex:row}
The shape $(3)$ admits exactly one standard tableau, so the expectation is
deterministic. Its hook lengths are $3, 2, 1$, with sum $6$. Then
\[
\frac{\sum h(c)}{n^2} \;=\; \frac{3+2+1}{9} \;=\; \frac{6}{9} \;=\; \frac{2}{3},
\qquad
\frac{n^2-1}{3} \;=\; \frac{9-1}{3} \;=\; \frac{8}{3},
\]
and $2/3 \ne 8/3$. The conjecture is false.
\end{counterexample}

\section{Alternative readings also collapse}

The statement is terse, so we check the plausible readings; all are false.

\paragraph{Square reading ($n \times n$ tableau, $n^2$ cells).}
The hook length of cell $(i,j)$ of an $n \times n$ square is
$(n-i)+(n-j)+1$, whence
\[
\sum h(c) \;=\; \sum_{i,j=1}^{n}\bigl((n-i)+(n-j)+1\bigr)
\;=\; n^2(n-1) + n^2 \;=\; n^3,
\qquad
\frac{\sum h}{n^2} = n .
\]
Setting $n = (n^2-1)/3$ gives $n^2 - 3n - 1 = 0$, whose discriminant $13$ is not
a perfect square; there is no integer solution, so the square reading fails for
\emph{every} $n \ge 1$. Example: the $3\times 3$ square has hook multiset
$\{5,4,4,3,3,3,2,2,1\}$, sum $27 = 3^3$, and $27/9 = 3 \ne 8/3$.

\paragraph{Divide-by-$n$ reading ($(\sum h)/n = (n^2-1)/3$, $n$ cells).}
Still $\sum h \le n^2$, so $(\sum h)/n \le n$, while $(n^2-1)/3 > n$ for every
$n \ge 4$ (as $n^2 - 3n - 1 > 0$ there). Small cases fail directly: for $n = 3$,
row shape gives $6/3 = 2 \ne 8/3$; for $n = 4$, row shape gives $10/4 = 5/2 \ne 5$.

\paragraph{Genuine randomization.}
The shape $(2,1)$ ($n = 3$ cells) has two standard tableaux, both with the hook
multiset $\{3,1,1\}$; the expectation is $E[\sum h] = 5$, and
$E[\sum h]/9 = 5/9 \ne 8/3$, $E[\sum h]/3 = 5/3 \ne 8/3$. Randomization over
tableaux of a fixed shape cannot move the hook-sum at all, and randomization
over shapes cannot exceed the cap of Fact~\ref{fact:cap}.

\section{Boundary}

\begin{itemize}
\item $n = 1$: LHS $= 1$ (either reading), conjectured value $0$. Already false.
\item $n = 2$: row shape gives $3/4$ (div.\ by $n^2$) or $3/2$ (div.\ by $n$)
  versus conjectured $1$; the $2\times 2$ square gives $8/4 = 2$ versus $1$.
  False.
\item $n \ge 3$: killed by the cap argument (div.\ by $n^2$), any shape;
  $n \ge 4$ killed by the cap under the divide-by-$n$ reading; the square
  reading dies for every $n$.
\end{itemize}

\section{Verification}

The repository contains:
\begin{itemize}
\item \texttt{reproduce.py} --- a standalone (stdlib-only) recomputation of every
  number above: hook multisets for shapes $(3)$, $(2,1)$, $(3,3,3)$, the
  expectation for shape $(2,1)$ via the hook length formula, the cap checks, and
  the discriminant argument.
\item \texttt{lean4/} --- a Lean~4 (v4.33.1, no Mathlib) development that
  concretizes the attack numbers as theorems proved by \texttt{rfl}/\texttt{decide}
  over \texttt{Nat}, including the counterexample $3+2+1 = 6$ with
  $6 \cdot 3 \ne 8 \cdot 9$ (i.e.\ $2/3 \ne 8/3$) and the overshoot
  $3 < n^2 - 1$ for $n \ge 3$. \texttt{Check.lean} audits every theorem with
  \texttt{\#print axioms}; all are reported as \emph{does not depend on any
  axioms} --- zero axioms, zero \texttt{sorry}.
\end{itemize}

\end{document}
